% !PRINTED 515
\weberpage{515}
%<<CONT
satisfied, and indeed the first or the second, according as
\weq{11}{\left(\frac{A}{m''}\right) = +1 \quad \text{or} \quad \left(\frac{A}{m''}\right) = -1}
holds.

2. \gesperrt{From this it follows that the sign} $\left(\dfrac{A}{m''}\right)=\pm 1$ \gesperrt{does not depend on the particular form} $(A,B,C)$\gesperrt{, but only on the class represented by this form, and is therefore a character of the class of forms.}
Now the invariant of a two-sided class is real, and the invariants of opposite classes are \gesperrt{conjugate imaginary.} But for two opposite forms $(A,B,C)$, $(A,-B,C)$ the sign (11) is the same, and from this it follows that $\Phi(\mu,u)$ has either real or conjugate imaginary roots, and consequently \gesperrt{real coefficients}, and therefore does not change when $\mu$ is replaced by the conjugate imaginary value $\mu'$. Now the two cases $m''\equiv 1$, $m''\equiv 3 \ (\mathrm{mod}\ 4)$ must be distinguished, because it depends upon this which term in (9) is imaginary. In the one case $i$ occurs only in the combination $i\sqrt{m'}$, in the other in $i\sqrt{m''}$, and accordingly we have:
\weq{12}{\begin{aligned}
\Phi(\mu,u) = \tfrac{1}{2}\bigl[\Phi(\mu,u)+\Phi(\mu',u)\bigr] &= \Psi(\sqrt{m''},u): m''\equiv 1 \\
&= \Psi(\sqrt{m'},u): m''\equiv 3
\end{aligned}
\ (\mathrm{mod}\ 4).}
The function $\Psi$, and consequently $\Phi$ also, has real coefficients and contains only one of the two square roots $\sqrt{m''}, \sqrt{m'}$. By our assumption $\sqrt{m''}$ is always \gesperrt{irrational}, while $\sqrt{m'}$ is rational only when $m'$ is a square. Since $H(u)$ has rational coefficients, it must be divisible, when $m''\equiv 1$, by $\Psi(\sqrt{m''},n)$ and $\Psi(-\sqrt{m''},n)$, and when $m''\equiv 3$ and $m'$ is not a square, by $\Psi(\sqrt{m'},n)$ and $\Psi(-\sqrt{m'},n)$.
% SOURCE-ORIGINAL-SIC: The registered page unambiguously prints n as the second argument in all four Psi terms, although equation (12) and H(u) indicate that u is mathematically expected.
Now $M$ changes its sign through the simultaneous changes of sign:
\weq{13}{
\begin{aligned}
(i,-i), (\sqrt{m''},-\sqrt{m''}): & m''\equiv 1 \\
(i,-i), (\sqrt{m'},-\sqrt{m'}): & m''\equiv 3
\end{aligned}
\ (\mathrm{mod}\ 4),}
and consequently, unless $m''\equiv 3$ and at the same time $m'$ is a square, we have in both cases:
\weq{14}{H(u) = \Phi(M,u)\,\Phi(-M,u),}
and this decomposition is impossible only when $m''\equiv 1 \ (\mathrm{mod}\ 4)$ and at the same time $m'$ is a square or $m'=1$.
% SOURCE-ORIGINAL-SIC: The registered page prints m'' congruent to 1 here, although the sentence introducing (14) and the preceding argument indicate the exceptional case m'' congruent to 3.
%CONT>>
